\lstinputlisting[language={[LaTeX]{TeX}},style=block,float,caption={Grundgerüst einer Bibliografie},label={lst:latex:literatur:template}]{code/latex_literatur_template.tex}

Die Erstellung einer Bibliografie erfolgt in mehreren Schritten.
Listing~\ref{lst:latex:literatur:template} zeigt den grundlegenden Aufbau.

\begin{compactenum}
\item Zunächst muss in der Präambel das Paket \texttt{biblatex} geladen werden.
\item Anschließend werden noch im Vorspann mit
      \lstinline[language={[LaTeX]{TeX}}]+\addbibresource+
      die Literaturdatenbanken geladen. Die Vorlage verwendet hierfür
      standardmäßig eine einzige Datei~\texttt{literatur.bib}. Für umfangreiche
      Arbeiten oder komplexe Literaturstudien ist es auch möglich, die
      Bibliografie aus mehreren \texttt{.bib}-Dateien zusammenzusetzen.
      Zweckmäßiger Weise legt man diese Dateien dann in einem separaten
      Unterverzeichnis ab. Jede Datei mit einem entsprechenden
      \lstinline[language={[LaTeX]{TeX}}]+\addbibresource+-Befehl
      geladen werden.
\item Die eigentliche Bibliografie wird innerhalb des Dokuments mit
      \lstinline[language={[LaTeX]{TeX}}]+\printbibliography+
      ausgegeben.
\end{compactenum}

Das Paket sowohl die oben genannten Befehle lassen sich durch eine Vielzahl von
Optionen angepassen. Hierfür sei auf die entsprechende Dokumentation in
\cite{pkgmanual_biblatex} verwiesen.

Die Quellenverweis ist mit dem
\lstinline[language={[LaTeX]{TeX}}]+\cite+-Befehl. Das Schlüsselwort für die
Fundstelle entspricht dem Schlüsselwort, welches in der \hologo{BibTeX}-Datei
vergeben wurde. Daneben gibt es noch eine Vielzahl weiterer Zitierbefehle, die
für verschiedene Anweungsfälle gedacht sind. Der folgende Absatz listet einige
Varianten beispielhaft auf. Der \LaTeX-Code hierzu ist in
Listing~\ref{lst:latex:literatur:cite} abgedruckt. Eine genaue Beschreibung
jeder Variante würde den Rahmen dieser Anleitung sprengen. Hierfür sei auf
\cite{pkgmanual_biblatex,rees_biblatex_cheat} verwiesen.

Diese Dokumentation ist unter \cite{urz_vorlage_manual}
veröffentlich. Der Quellenverweis kann auch per Fußnote
erfolgen.\footcite{urz_vorlage_manual} Man kann auch explizit
nur die Autoren zitieren: \citeauthor{urz_vorlage_manual} oder
auch lediglich den Titel: \citetitle{urz_vorlage_manual}.
Selbst komplette Quellenangaben im Fließtext sind möglich:
\fullcite{urz_vorlage_manual} oder auch in der
Fußnote.\footfullcite{urz_vorlage_manual}


\lstinputlisting[language={[LaTeX]{TeX}},style=block,float,caption={Literaturverweise},label={lst:latex:literatur:cite},morekeywords={footcite,citeauthor,citetitle,fullcite,footfullcite}]{code/latex_literatur_cite.tex}
